\section{课程感想}

\begin{frame}{会用与精通之间}
\begin{itemize}
  \item 会几条命令、能装个软件，和真正精通 Linux 之间，差距比想象中大得多
  \item 比如课程中学了 \texttt{tr}，看起来只是字符替换的小工具——但它支持 POSIX 字符类、正则组合、管道串联，和 \texttt{grep}、\texttt{sort}、\texttt{uniq} 配合起来可以替代很多几十行的脚本
  \item 又比如 Vim：会 \texttt{h/j/k/l} 算是入门，但宏录制、文本对象、寄存器操作——这些真正提升效率的功能，需要大量刻意练习
  \item Linux 的每个工具都像冰山，表面看起来简单，水下的深度需要不断探索和积累
  \item 好的工具真的很多——\texttt{jq}、\texttt{tr}、\texttt{grep}、\texttt{sed}、\texttt{awk}——但如何把它们组合成高效的 workflow，不是学一两个命令就能做到的，需要\textbf{持续打磨}
\end{itemize}
\end{frame}

\begin{frame}{AI 工具的实际帮助}
\begin{itemize}
  \item 这门课的学习过程中，AI 工具带来了很大的便利
  \item 写 Shell 脚本时，遇到不熟悉的语法就丢给 Claude Code 或 opencode 解释，几乎立即能得到带示例的说明——比翻 man 手册高效太多
  \item 在写 cal.sh 计算 PI 时，bc 的精度参数和 Ramanujan 公式的实现细节，AI 可以帮忙快速验证逻辑，减少了很多试错成本
  \item opencode 尤其在 LaTeX 部分发挥了作用：修改 Beamer 主题、调整排版、写 Makefile——AI 协助下从零到完成一份 PDF 相当顺畅
  \item openclaw 等工具在处理批量文件操作时节省了大量手工劳动，比如批量重命名、格式化代码
\end{itemize}
\end{frame}

\begin{frame}{AI 时代的思考方式}
\begin{itemize}
  \item 但更重要的是 AI 没有替代思考：最终的理解、判断、决策——比如什么时候用管道什么时候写脚本、哪种排版更合适——还是需要自己来做
  \item AI 给出方案后，仍然需要人去评估是否正确、是否适合当前场景——批判性思维比任何时候都重要
  \item 未来这些工具只会越来越强，持续尝试新事物、更新自己的工具链，应该是每个技术从业者的基本素养
  \item 一个很深的体会是：AI 时代，知道问什么比知道怎么做一个技术细节更重要——方向对了，AI 能帮你跑得更快
\end{itemize}
\end{frame}
